% Thesis Proposal in LaTeX (English)
\documentclass[12pt]{article}
\usepackage{amsmath, amssymb, graphicx, physics, bm, geometry}
\geometry{margin=1in}

\title{\textbf{Thesis Proposal: A Reduced-Order Aerodynamic State Modeling Framework for Real-Time Prediction and Control Using Sparse Sensors}}
\author{PhD Candidate: [Your Name] \\ Supervisor: [Supervisor Name] \\ Institution: [Institution]}
\date{\today}

\begin{document}
\maketitle

\section{Motivation}
Accurate aerodynamic characterization of bluff bodies---such as cylinders, simplified mannequins, and human athletes---is essential for performance optimization and control. Computational Fluid Dynamics (CFD) provides high-fidelity predictions for fixed geometries and controlled conditions, yet it remains computationally prohibitive for real-time applications, dynamic postures, and feedback-driven control. Furthermore, CFD cannot be directly integrated with noisy experimental sensor data in operational settings.

To address these gaps, this thesis proposes a unified, physics-informed Reduced-Order Modeling (ROM) and data-assimilation framework capable of predicting aerodynamic forces in real time using sparse pressure measurements. The goal is to develop a systematic modeling architecture applicable across a sequence of increasingly complex cases, from canonical flows to human posture aerodynamics.

\section{Research Question}
Can a coherent-structure-based reduced-order aerodynamic state model, equipped with a dynamic predictor and sensor-based corrector, be constructed to estimate and predict aerodynamic forces of bluff bodies in real time, and can this framework be extended across cases of increasing geometric and flow complexity?

\section{Hypotheses}
\begin{itemize}
    \item \textbf{H1 (Coherent Structures):} Bluff-body wakes exhibit dominant coherent structures that can be captured in a low-dimensional modal or latent space.
    \item \textbf{H2 (Reduced Dynamics):} The temporal evolution of these structures can be represented by a reduced-order dynamic system preserving the relationship between state and aerodynamic forces.
    \item \textbf{H3 (Sensor-State Mapping):} Sparse pressure sensors contain sufficient information to estimate the reduced aerodynamic state and reconstruct drag in real time.
    \item \textbf{H4 (Framework Extensibility):} A common modeling architecture---state representation, predictor, and corrector---can be adapted to progressively complex cases with minimal structural modifications.
\end{itemize}

\section{Objectives}
\subsection*{Primary Objectives}
\begin{enumerate}
    \item Develop a reduced aerodynamic state representation using modal decomposition (POD/SPOD) or nonlinear latent embeddings.
    \item Construct a dynamic predictor (e.g., Galerkin ROM, operator inference, Koopman model, or neural ODE) governing state evolution.
    \item Design a measurement operator and data-assimilation corrector (EKF/EnKF/variational methods) to estimate aerodynamic state from sparse sensors.
    \item Develop a data-driven state map (clustering + transition model) as an alternative predictor for drag reconstruction.
    \item Validate extensibility across three stages: cylinder, simplified 3D dummy, and human body with limited degrees of freedom.
\end{enumerate}

\section{Methodology}
\subsection{Phase 1: Cylinder Flow (Proof of Concept)}
\begin{itemize}
    \item Acquire experimental data for three flow velocities using four pressure sensors.
    \item Extract coherent structures (POD/SPOD) and construct low-dimensional states.
    \item Train ROM and clustering-based state-transition models.
    \item Evaluate drag reconstruction accuracy and cross-velocity generalization.
\end{itemize}

\subsection{Phase 2: Simplified 3D Dummy}
\begin{itemize}
    \item Instrument dummy with sparse sensors at aerodynamically relevant locations.
    \item Construct geometry-specific state representation and dynamic model.
    \item Validate relationships between posture variation, state signatures, and aerodynamic forces.
\end{itemize}

\subsection{Phase 3: Human Body with Restricted Posture Space}
\begin{itemize}
    \item Define a limited set of posture degrees of freedom.
    \item Build a reduced-order predictive model for drag and lateral forces.
    \item Demonstrate potential for posture optimisation toward reduced drag.
\end{itemize}

\section{Expected Contributions}
\begin{itemize}
    \item A general aerodynamic state modeling framework based on coherent structures.
    \item A real-time predictor--corrector system for estimating drag using sparse sensors.
    \item A data-driven aerodynamic state map applicable as an alternative to physics-based ROM.
    \item Multi-level validation from canonical to biomechanically relevant scenarios.
    \item Foundations for real-time aerodynamic optimisation for sports and human-in-the-loop control.
\end{itemize}

\section{Timeline}
\begin{enumerate}
    \item Year 1: \begin{itemize}
                        \item Cylinder experiments, ROM design, sensor-state mapping.
                        \item Dummy experiments, model extension, data assimilation.
                    \end{itemize}
    \item Year 2: Human posture modelling, optimisation layer, thesis consolidation.
\end{enumerate}

\end{document}
